\chapter{Grassmannians and the space of lines in \texorpdfstring{$\mathbb{P}^3$}{P3}}
\label{grassmannians}

In this chapter we will introduce Grassmann varieties, or Grassmannians, and focus on $\mathbb{G}(1, 3)$, which is the space of lines in $\mathbb{P}^3$. We will also introduce the Plücker embedding, which is a way to embed the Grassmannian into a projective space, giving it the structure of a projective variety. For a more detailed introduction to Grassmannians, we refer to (Eisenbud, Harris)(!?).

A Grassmannian is a projective variety whose closed points correspond to linear subspaces of a certain dimension in an $n$-dimensional vector space. As a set, $Gr(k, V)$ is the set of all $k$-dimensional linear subspaces of the vector space $V$. In our case, $V$ will be a vector space over $\mathbb{C}$. A $k$-dimensional linear subspace of an $n$-dimensional vector space $V$ is the same as a $(k - 1)$-dimensional linear subspace of $\mathbb{P}V \cong \mathbb{P}^{n-1}$. 

To give it the structure of a projective variety, we give an inclusion in a projective space, called the Plücker embedding, and showing that the image is the zero locus of a certain collection of homogeneous polynomials.

\section{The Plücker embedding}

